\begin{resumo}[Abstract]
\begin{otherlanguage*}{english}
This work develops a computational routine to extract, organize, and analyze
voltage maxima from transmission line energization simulations in the
\textit{Alternative Transients Program} (ATP). The evaluation uses two
scenarios, with and without a pre-insertion resistor, five simulation sample
sizes, three terminals, and three phases, totaling 204,300 scalar
observations. The focus is a comparative analysis of K-Means, DBSCAN, and
hierarchical upper-tail selection using Silhouette and Davies--Bouldin
indices, cluster counts, and applicability status. The comparison examines
how results vary with the number of simulations, without interpreting
geometric separation as proof of numerical error. The probabilistic
analysis uses empirical frequencies and confidence intervals; rejection of
the normality hypothesis in 88 of the 90 groups prevents adopting a common
Gaussian distribution as a reference for validating candidates. The routine
preserves traceability and produces reproducible evidence for further
investigation. Causal validation by an expert, using waveforms and repeated
simulations, remains future work outside the scope of this thesis.

\vspace{\onelineskip}
\noindent
\textbf{Keywords}: Transient overvoltages. Anomaly detection. Machine learning.
ATP/EMTP. Probabilistic analysis.
\end{otherlanguage*}
\end{resumo}
